\documentclass[12pt]{article}
\usepackage[utf8]{inputenc}
\usepackage[T1]{fontenc}
\usepackage{hyperref}
\title{Sya Language Resources for K-12 Curriculum Development}
\author{Compiled by Claude Sonnet 5.5 (claude-sonnet-5-5) for Andr\'as Kornai}
\date{}
\begin{document}
\maketitle

\section{Introduction}
Sya (also spelled Sha; ISO 639-3 scw) is a Chadic language of the Afro-Asiatic family. It is spoken in Sha district, Bokkos Local Government Area, Plateau State, Nigeria, on the western edge of the Jos Plateau. Speaker estimates conflict: about 1,000 (Wikipedia, 2018), about 3,000 (a 1998 figure), about 6,700 (Joshua Project people-group population) and about 7,000 (a numerals database). The closely related Mundat language is spoken in a village in the same district. Searches found no school materials and no standard orthography. Resources for this language are very scarce.

\section{Primary Educational Resources}
No primers, textbooks or school curricula in Sya turned up in any search. Wikipedia states that the language is used for face-to-face communication and has no standardized written form.
\begin{itemize}
\item Overview of the language: \url{https://en.wikipedia.org/wiki/Sha_language}
\end{itemize}

\section{Linguistic and Reference Materials}
\begin{itemize}
\item Sya numerals page (Max Planck Institute channumerals database, data supplied by Mark Rosenfelder in October 2023): \url{https://lingweb.eva.mpg.de/channumerals/Sya.htm}
\item ASJP word-list entry for Saya/Sya: \url{https://asjp.clld.org/languages/SAYA}
\item Newman, Chadic Classification and Index (2013), for the language's place in Chadic: \url{https://megatchad.cepam.cnrs.fr/publications/Newman-2013-Chadic-Classification-and-Index.pdf}
\end{itemize}
I did not find a published grammar, dictionary or orthography statement.

\section{Religious and Community Resources}
Joshua Project lists scripture status as ``Translation Started'' and notes that audio recordings exist. I did not open or verify any scripture text.
\begin{itemize}
\item Joshua Project, Sha in Nigeria: \url{https://joshuaproject.net/people_groups/14882}
\end{itemize}

\section{Digital and Audio Resources}
No digital tools, dictionaries or audio corpora were found. Joshua Project mentions recordings, but I did not see them directly.

\section{Curriculum Development Context}
There is no known schooling in Sya. Any curriculum would first need an agreed orthography, a documented lexicon and grammar, and community partnership. These do not appear to exist in public sources today. The first steps would be to contact SIL and Nigerian linguistic institutions and local Bokkos community leaders. I found no named organisation or contact for this, so none is listed here.
\end{document}